% !Mode:: "TeX:UTF-8"
%!TEX program  = xelatex

\documentclass[bwprint,fontset=fandol]{gmcmthesis}

\hypersetup{hidelinks}

%========================================================
% 题目、摘要、关键词标题：隶书，小二
% 字体文件位置：Fonts/SIMLI.TTF
%========================================================

% 定义真正的隶书字体
\setCJKfamilyfont{mylishu}[
    Path = Fonts/
]{SIMLI.TTF}

% 将模板中的 \lishu 指向上传的真正隶书
\renewcommand{\lishu}{\CJKfamily{mylishu}}

%========================================================
% 将“关键词”三个字统一修改为：隶书、小二
%========================================================
\makeatletter

\renewcommand{\keywords}[1]{%
    \renewcommand{\gmcm@tokens@keywords}{#1}%
    \par
    \vskip1ex
    {\noindent\zihao{-2}\lishu
    \gmcm@cap@keywordsname：}%
    ~{\gmcm@tokens@keywords}%
}

\makeatother
\usepackage{graphicx}   % 必须，否则 \includegraphics 报错

\usepackage[framemethod=TikZ]{mdframed}

%\usepackage{fontspec}  % for Consolas & Courier New
%\lstset{basicstyle=\small\fontspec{Consolas}}
%\lstset{basicstyle=\small\fontspec{Courier New}}
%\setmonofont{Consolas}
%\setcounter{tocdepth}{3}   % 调整目录深度
\usepackage{siunitx}
\usepackage{colortbl}
\definecolor{color1}{rgb}{0.78,0.88,0.99}
\definecolor{color2}{rgb}{0.36,0.62,0.84}
%\definecolor{color3}{rgb}{0.8235,0.8706,0.9373}
\definecolor{color3}{rgb}{0.88,0.92,0.96}
\definecolor{color4}{rgb}{0.96,0.97,0.98}%{0.9176,0.9373,0.9686}

% 算法
\usepackage[noend]{algpseudocode}
\usepackage{algorithmicx,algorithm}
\floatname{algorithm}{算法}
\renewcommand{\algorithmicrequire}{\textbf{输入:}}
\renewcommand{\algorithmicensure}{\textbf{输出:}}

%\numberwithin{equation}{section}
%\numberwithin{figure}{section}
%\numberwithin{table}{section}


\newcommand{\red}[1]{\textcolor{red}{#1}}
\newcommand{\blue}[1]{\textcolor{blue}{#1}}


%===================== 建模论文题目 ======================

\title{山区洪涝灾害下无人机运输与通信协同优化
}
\baominghao{\hspace{6em} 26106350114} %参赛队号
\schoolname{\hspace{5.8em} 西南大学} %学校名称
\membera{\hspace{6em} 谢文哲} %队员A
\memberb{\hspace{6em} 章林怡} %队员B
\memberc{\hspace{6em} 刘宇豪} %队员C

%=======================================================


\begin{document}

%生成标题
\maketitle

%填写摘要
\begin{abstract}

山区洪涝灾害往往伴随着\textbf{道路阻断、通信受损与能源受限}等多重约束，使传统地面运输难以及时完成应急物资配送。无人机具有机动性强、受地面道路影响小等优势，但其实际救援能力同时受到载荷、地形、能耗、配送时限、通信覆盖以及设备周转等因素制约。针对山区洪涝灾害场景下无人机应急运输与通信保障的协同优化问题，本文结合服务区空间位置、数字高程模型（DEM）、货箱需求、异构无人机参数、电池资源及无线通信链路参数，建立了由\textbf{单点运输能力分析、多机多架次调度、运输--通信联合调度和任务分区资源配置}构成的递进式优化模型，并从运输效率、任务及时性、能源消耗和资源占用等多个维度对方案进行评价。

针对问题一，为准确刻画不同机型在山区地形条件下的运输能力，首先利用\textbf{DEM栅格采样}提取调度中心与各服务区直线路径上的最高地形高程，并按照“爬升--巡航--下降”的飞行过程建立时间与能耗计算模型。在载荷影响等效航程且返航需保留安全能量余量的条件下，利用\textbf{二分搜索}求解各“机型--服务区”组合的最大安全载荷。在此基础上，将不可拆货箱的组批过程转化为同时考虑载质量、装载体积与安全能量约束的\textbf{多维装箱优化问题}，以架次数、总运输能耗和累计作业时间构造分层多目标优化模型，获得各服务区的货箱组批与机型选择方案。同时改变返航安全余量进行敏感性分析，刻画安全冗余增加对最大可运输载荷及架次数的影响规律。

针对问题二，在单架次可访问多个服务区的条件下，进一步考虑\textbf{异构无人机、共享电池、配送时限与多架次周转}之间的耦合关系。本文构建带容量、能量和时间窗约束的异构多旅行车辆路径模型，将货箱组批、服务区访问顺序、运输机型、实体无人机、电池资源及架次开始时刻进行联合决策。为降低大规模组合搜索的复杂度，采用\textbf{“可行架次生成--资源调度优化”两阶段求解框架}：首先生成满足载荷、体积和能量限制的候选运输架次，再结合医疗物资期望送达时间和首批保障截止时间，通过自适应邻域搜索与整数优化对路线及资源时序进行协调。模型能够输出逐架次运输路线、逐箱送达时刻以及无人机和电池占用情况，并通过任务时序和资源冲突检验保证方案的实际可执行性。

针对问题三，在运输调度基础上进一步引入山区地形遮挡和无线传播约束。首先依据运输无人机的连续飞行轨迹，结合DEM完成\textbf{视距（LOS）遮挡判定}，并根据双向链路预算计算运输无人机与固定网关之间的实时通信可用性。对于直连链路失效区段，在DEM覆盖区域内生成满足悬停高度约束的候选中继位置，并预计算“运输无人机--中继无人机--固定网关”两段链路的可服务关系，从而将连续空间中继部署问题转化为有限候选位置上的\textbf{时空覆盖问题}。进一步建立运输无人机与中继无人机联合调度模型，同时优化运输架次、中继悬停位置、服务时段及能源组件配置，在保证运输无人机\textbf{全飞行过程连续通信}的前提下，综合降低物资延误、联合任务完成时间、运输与中继总能耗以及设备使用强度，并通过时间离散精度的敏感性分析验证通信判定结果的稳定性。

针对问题四，为研究灾后分区独立救援条件下的资源配置问题，根据问题三中已经确定的运输架次和通信保障关系构建\textbf{服务区任务耦合图}。若多个服务区位于同一运输架次或存在不可分割的联合保障关系，则将其合并为不可拆分的任务单元，并分别建立两分区和三分区优化模型。在保持原运输与中继任务安排不变的条件下，以\textbf{资源配置规模、资源冗余、组间工作量均衡和现有库存缺口}为主要评价指标进行分区优化；同时基于任务占用时间区间计算各分区运输无人机、共享电池、中继无人机及能源组件的最大并发需求，从而得到各分区独立执行所需的最小资源配置，并识别潜在资源瓶颈及其形成原因。

综上，本文将\textbf{地形约束、运输调度、能源周转与连续通信保障}统一纳入山区应急无人机救援决策框架，在保证货箱不可拆分、任务时限、返航安全以及通信连续性的基础上，实现了由单架次能力分析到多机协同、再到运输--通信联合优化与灾区资源分区配置的逐层扩展。所建立模型具有明确的物理含义和较好的可解释性，可为山区洪涝灾害条件下无人机应急运输与通信资源协同配置提供定量化决策依据。

\keywords{山区洪涝灾害\quad  无人机应急运输\quad 异构多架次调度\quad  通信中继  \quad  协同优化}
\end{abstract}

\pagestyle{plain}

%目录 不推荐加
\maketoc

\clearpage

\section{问题重述}

\subsection{引言}   % 问题的背景
山区洪涝灾害具有突发性强、地形复杂和基础设施易受损等特点。持续强降雨可能引发山洪、滑坡及道路塌方，并进一步造成道路、供电和通信设施同时中断，使部分受灾区域在短时间内难以通过传统地面交通获得物资补给。题目所描述的广西横州市镇龙乡洪涝灾害中，就曾出现断路、断电、断网以及大量群众受困的情况，食品、饮用水、药品及抢修设备需要依靠徒步或空中方式运入灾区。 因此，利用受地面道路条件影响较小的无人机开展应急物资运输，对提高山区灾害救援效率具有重要意义。

然而，山区无人机救援并非简单的物资配送问题。一方面，复杂地形会影响无人机的飞行高度、航程及能源消耗，不同货箱的重量、体积和配送时限也使货箱组批、路径规划、机型选择和电池周转相互制约；另一方面，山体遮挡可能导致运输无人机与固定网关之间通信中断，需要利用中继无人机提供持续通信保障。由此，运输任务、通信保障和有限装备资源之间形成了明显的耦合关系。题目也明确要求同时考虑载荷、地形净空、能源储备、任务时限及通信链路等多种因素。

基于此，本文以题目给出的调度中心、15个服务区、80个不可拆货箱以及运输无人机、中继无人机、共享电池和DEM等数据为基础，围绕单点运输能力、多机多架次调度、运输与通信协同以及任务分区资源配置四个层次展开建模，在保证运输安全、配送时效和通信连续性的前提下，对有限无人机及能源资源进行合理配置，为山区洪涝灾害下的无人机应急救援提供定量化决策依据。
\subsection{问题的提出}

%\subsubsection{问题的提出内容一}

\noindent 根据题目给定的调度中心、服务区、货箱需求、DEM地形数据、运输与中继无人机参数、电池资源及通信链路参数，需要依次解决以下四个相互递进的问题。

\textbf{问题一：} 单点往返运输能力与货箱组批。针对不同机型在各服务区执行
$O_{01}\rightarrow S_i\rightarrow O_{01}$
往返任务的场景，结合DEM地形、载荷和能耗约束，确定各机型的\textbf{最大安全载荷}。在货箱不可拆分的条件下，优化\textbf{货箱组批方案}，并分析返航安全余量变化对运输能力和组批结果的影响。

\textbf{问题二：} 异构无人机多点多架次运输调度。在允许单架次访问多个服务区的条件下，综合考虑载荷、配送时限、无人机数量及共享电池周转等约束，联合优化\textbf{货箱组批、访问顺序、无人机分配和架次时序}，以提高配送及时性并降低任务完成时间、能耗和架次数。

\textbf{问题三：} 通信约束下的运输与中继联合调度。在问题二的基础上引入山区地形遮挡和通信链路约束，当运输无人机与固定网关 $G_{01}$ 无法直连时，通过中继无人机提供通信保障。联合优化\textbf{运输任务与中继无人机的悬停位置、服务时段及能源配置}，保证运输全过程通信连续，并兼顾时效与能耗。

\textbf{问题四：} 救援任务分区与资源配置。在保持问题三运输与通信任务关系不变的条件下，将 $15$ 个服务区分别划分为 $2$ 个和 $3$ 个任务组。比较不同分区方案下的\textbf{资源需求、资源冗余、工作量均衡及库存缺口}，确定更合理的分区与资源配置方案。

\clearpage

\section{模型假设与符号说明}
\subsection{模型假设}
\subsection{符号说明}
\begin{table}[htp!]
\centering
\renewcommand\arraystretch{1.2} %定义表格高度
\newcolumntype{L}{>{\quad}X}
\newcolumntype{P}[1]{>{\centering\arraybackslash}p{#1}}
\newcolumntype{C}{>{\centering \arraybackslash}X}
\newcolumntype{R}{>{\raggedright \arraybackslash}X}
\begin{tabularx}{0.9\textwidth}{|P{1.8cm}|L|}
  \hline %\toprule
  符号    &   \quad 意义 \\
  \hline %\midrule
  $ O $  &  临时调度中心    \\
  \hline
  $ S_i $  &  第 $ i $ 个服务区，$ i=1,2,\ldots,15 $    \\
  \hline
  $ g $  &  运输无人机机型，$ g \in \{A,B,C\} $    \\
  \hline
  $ m_g $  &  机型 $ g $ 的无人机空机质量/kg    \\
  \hline
  $ Q_g $  &  机型 $ g $ 的额定最大载货质量/kg    \\
  \hline
  $ V_g $  &  机型 $ g $ 的最大装载体积/m$^3$    \\
  \hline
  $ E_g $  &  机型 $ g $ 单组电池可用能量/kWh    \\
  \hline
  $ p_g $  &  机型 $ g $ 的返航安全余量比例    \\
  \hline
  $ v_g $  &  机型 $ g $ 的水平巡航速度/(m·s$^{-1}$)    \\
  \hline
  $ v_g^\uparrow, v_g^\downarrow $  &  机型 $ g $ 的爬升速度和下降速度/(m·s$^{-1}$)    \\
  \hline
  $ \eta_g $  &  爬升附加能耗效率    \\
  \hline
  $ R_g^0, R_g^F $  &  机型 $ g $ 的空载、满载标准航程/m    \\
  \hline
  $ d_{ij} $  &  节点 $ i $ 到节点 $ j $ 的水平测地距离/m    \\
  \hline
  $ H_{ij} $  &  航线 $ i \to j $ 的计划巡航海拔/m    \\
  \hline
  $ h_{ij}^\uparrow, h_{ij}^\downarrow $  &  航线 $ i \to j $ 的爬升高度、下降高度/m    \\
  \hline
  $ q $  &  当前有效载荷/kg    \\
  \hline
  $ R_g(q) $  &  载荷 $ q $ 下的等效航程/m    \\
  \hline
  $ E_{ij}(g|q) $  &  机型 $ g $ 在载荷 $ q $ 下执行航线 $ i \to j $ 的运输能耗/kWh    \\
  \hline
  $ t_{ij}(g) $  &  机型 $ g $ 执行航线 $ i \to j $ 的飞行时间/s    \\
  \hline
  $ B_i $  &  服务区 $ S_i $ 的货箱集合    \\
  \hline
  $ m_b, v_b $  &  货箱 $ b $ 的质量/kg 与体积/m$^3$    \\
  \hline
  $ q_g^{\mathrm{safe}} $  &  机型 $ g $ 服务 $ S_j $ 时的最大安全载荷/kg    \\
  \hline
  $ p $  &  某一可行班驶模式    \\
  \hline
  $ c_p $  &  班驶模式 $ p $ 对各类货箱的数量向量    \\
  \hline
  $ N $  &  全部服务区总运输频次    \\
  \hline
  $ E $  &  全部班次总运输能耗/kWh    \\
  \hline
  $ T $  &  全部班次累计作业时间/s    \\
  \hline
\end{tabularx}
\end{table}


%\begin{tabularx}{\textwidth-18pt}{XXX}
%\hline
%Input & Output& Action return \\
%\hline
%DNF &  simulation & jsp\\
%\hline
%\end{tabularx}


\section{问题一的建模与求解}

\subsection{问题一分析}
问题一要求在不考虑实体无人机和共享电池调度的条件下，研究“临时调度中心 $O01$---单个服务区 $S_i$---$O01$”的直接往返运输能力与货箱组批问题。由于每一架次只能服务一个服务区，15 个服务区之间不存在跨区组批，因此该问题可以先对每个“机型—服务区”组合建立独立的飞行时间与能耗模型，再在各服务区内部求解不可拆货箱的最优组批，最后汇总得到全局结果。
与普通装箱问题相比，本题的关键难点在于：第一，载荷上限不仅由机型额定载质量决定，还受 DEM 地形引起的爬升能耗、航程衰减和返航安全余量共同约束；第二，货箱不可拆分，同时受到质量和体积双重限制；第三，题目要求综合考虑往返架次数、运输能耗与累计作业时间，三者之间存在明显的层级权衡。因此，本文采用“物理可行性计算—可行组批模式枚举—状态压缩动态规划”的思路进行精确求解，整体流程如图~\ref{fig:problem1_flow}所示。
\begin{figure}[htbp!]
  \centering
  \fbox{%
    \begin{minipage}[c][5cm][c]{0.85\textwidth}
      \small（占位框，待替换为真实流程图）
    \end{minipage}%
  }
  \caption{问题一求解流程图}
  \label{fig:problem1_flow}
\end{figure}
\subsection{基本约定}
（1）每一运输架次均采用 $O01 \to S_i \to O01$ 的直接往返形式，只服务一个服务区；同一服务区允许使用多个架次完成配送，但不得跨服务区组批。

（2）任意两个任务节点之间的水平航迹取两节点连线；巡航海拔由该连线经过的 $30\,\mathrm{m}$ DEM 像元最高地形高程加 $50\,\mathrm{m}$ 安全净空确定。

（3）$O01$ 的作业高度取地面海拔，服务区作业高度取地面海拔以上 $30\,\mathrm{m}$；每个航段均按"爬升—水平巡航—下降"三个阶段计算飞行时间。

（4）下降阶段不单独计入附加势能消耗；运输总能耗由载荷相关水平航程能耗与爬升附加能耗构成。

（5）问题一不考虑实体无人机数量、共享电池数量及其充电周转，只考察单架次物理可行性与货箱组批；货箱不可拆分且每个货箱只能被安排一次。
\subsection{航段时间与能耗模型}

\subsection{算法示例}

数学建模求解算法示例：
\begin{center}
\begin{minipage}{0.8\textwidth}
\begin{algorithm}[H]%[!htp]
\caption{算法的名字} %算法的名字
{\bf 输入:} %算法的输入， \hspace*{0.02in}用来控制位置，同时利用 \\ 进行换行
input parameters A, B, C\\
{\bf 输出:} %算法的结果输出
output result
\begin{algorithmic}[1]
\State some description 算法介绍 % \State 后写一般语句
\For{condition} % For 语句，需要和EndFor对应
  \State ...
  \If{condition} % If 语句，需要和EndIf对应
    \State ...
  \Else
    \State ...
  \EndIf
\EndFor
\While{condition} % While语句，需要和EndWhile对应
  \State ...
\EndWhile
\State \Return result
\end{algorithmic}
\end{algorithm}
\end{minipage}
\end{center}
\vspace{2ex}

\subsection{表格}

使用 tabular 环境, 如下表格: 表~\ref{tab:foo}. 通过 \verb|autoref| 引用表格: \autoref{tab:foo}.

这里可使用命令 \verb|\tabcolsep| 和 \verb|\arraystretch| 分别控制列间距和行间距.

\begin{table}[htp!]
\centering
\setlength{\tabcolsep}{12pt}  % 6pt standard
\renewcommand{\arraystretch}{1.2}
\caption{学术活动安排样例}
\label{tab:foo}
\begin{tabular}{|c|c|c|c|}
\hline
\textbf{日期}  & \textbf{地点} & \textbf{活动名称} & \textbf{备注} \\ \hline
2024年8月1日      & 上海       & 学术研讨会      & 主题：人工智能 \\ \hline
2024年8月15日    & 北京       & 学术交流会      & 重点：数学建模 \\ \hline
2024年9月1日      & 深圳       & 研究研讨会      & 主题：数据科学 \\ \hline
2024年10月15日  & 广州       & 创新论坛         & 重点：科技创新 \\ \hline
\end{tabular}
\end{table}


\clearpage

三线表
\begin{table}[htp!]
\newcolumntype{L}{X}
\newcolumntype{C}{>{\centering \arraybackslash}X}
\newcolumntype{R}{>{\raggedright \arraybackslash}X}
\centering
\caption{某校学生升高体重样本}
\label{tab:heightweight}
\begin{tabularx}{0.9\textwidth}{CCCC}
\toprule
序号 & 年龄 & 身高 & 体重 \\
\midrule
001 & 15 & 156 & 42 \\
002 & 16 & 158 & 45 \\
003 & 14 & 162 & 48 \\
004 & 15 & 163 & 50 \\
\cmidrule{2-4}
平均 & 15 & 159.75 & 46.25 \\
\bottomrule
\end{tabularx}
\end{table}

% 某行业产量与生产费用的数据
% \begin{table}[htp!]
% \centering
% \caption{某行业产量与生产费用的数据}%\label{}
% \newcolumntype{Y}{>{\centering\arraybackslash}X}
% \newcolumntype{Z}{!{\vline}@{\color{white}\vrule width \doublerulesep}!{\vrule}}%自定义列格式(双线)
% \begin{tabularx}{0.94\textwidth}{c|c|YZc|c|Y}
% \Xhline{0.9pt}
% 企业编号&  产量(台)&生产费用(万元)&企业编号&产量(台)&生产费用(万元)\\\Xcline{1-3}{0.6pt}\Xcline{4-6}{0.6pt}
% 1 & 40 & 130 & 7&  84&  165\\
% 2 & 42 & 150 & 8&  100&  170\\
% 3 & 50 & 155 & 9&  116&  167\\
% 4 & 55 & 140 & 10&  125&  180\\
% 5 & 65 & 150 & 11&  130&  175\\
% 6 & 78 & 154 & 12&  140&  185\\\Xhline{0.72pt}
% \end{tabularx}
% \end{table}

研究生数学建模2019年F题结果示例
\begin{table}[htp!]
\centering
\caption{问题1结果1 (左) 与 问题2结果 (右)}
\begin{minipage}[h]{0.48\linewidth}
\renewcommand\arraystretch{1.2} %定义表格高度
\newcolumntype{Y}{>{\centering\arraybackslash}X}
\begin{tabularx}{0.9\textwidth}{|Y|Y|Y|}
  \hline
  数据集1  &  数据集1 & 数据集2  \\
 \hline
  A问题1   & A问题1 &  A问题1     \\
  \hline
  503   &  503     & 163     \\
  294   &  200     & 114      \\
  91    &  80      & 8     \\
  607   &  237     & 309      \\
  540   &  170     & 305    \\
  250   &  278     & 123    \\
  340   &  369     & 45      \\
  277   &  214     & 160    \\
  B     &  397     & 92    \\
        &  B       & 93    \\
        &          & 61        \\
        &          & 292       \\
        &          & B         \\
104861  & 103518   & 109342 \\
\hline
\end{tabularx}
\end{minipage}
\begin{minipage}[h]{0.48\linewidth}
\renewcommand\arraystretch{1.2} %定义表格高度
\newcolumntype{Y}{>{\centering\arraybackslash}X}
\begin{tabularx}{0.9\textwidth}{|Y|Y|Y|}
  \hline
  数据集1  &  数据集1 & 数据集2  \\
  \hline
  A问题2  &  A问题2  & A问题2    \\
  \hline
   503    & 503     & 163  \\
  294    & 200     & 114   \\
   91     & 80     & 8   \\
   607    & 237     & 309   \\
   540    & 170    & 305  \\
   250    & 278    & 123  \\
   340    & 369     & 45   \\
   277    & 214    & 160  \\
  B      & 397   & 92 \\
        &  B     &  93      \\
         &         &  61      \\
        &         &   292     \\
        &         &   B     \\
 104917  &  103563  &109427 \\
\hline
\end{tabularx}
\end{minipage}
\end{table}

\clearpage
\begin{table}[htp!]
%\small
\centering
\renewcommand\arraystretch{1.2} %定义表格高度
\newcolumntype{Y}{>{\centering\arraybackslash}X}
\caption{问题3结果}
\begin{tabularx}{0.9\textwidth}{|Y|Y|Y|Y|Y|Y|}
  \hline %定义表格宽度
  数据集1  &  数据集1  &  \multicolumn{2}{c|}{数据集2 (无问题点)}  & \multicolumn{2}{c|}{数据集2 (有问题点)} \\
  \hline
  A问题3   & A问题3   &  \multicolumn{2}{c|}{A问题3}    &  \multicolumn{2}{c|}{A问题3}      \\
  \hline
  503      & 503     & 169      & 73      & 169     & 73   \\
  \hline
  69       & 69      & 322      & 249     & 322     & 249   \\
  \hline
  506      & 506      & 270     & 274     & 270     & 274   \\
  \hline
  371      & 371      & 89      & 12      & 89      & 12   \\
  \hline
  183      & 183      & 236     & 216     & 236     & 216  \\
  \hline
  194      & 194      & 132     & 16      & 132     & 16   \\
  \hline
  450      & 450      & 53      & 282     & 53      & 282   \\
  \hline
  286      & 113      & 112     & 84      & 112     & 141  \\
  \hline
  485      & 485      &  268    & 287     & 268     & 291 \\
  \hline
 \red{B (9D)}~~     & 248      & 250     & 99      & 250     &161 \\
 \hline
   & \red{B (10D)}   & 243     & \red{B (21D)}   & 243    & \red{B (21D)} \\
   \hline
    &          &         &         &         &       \\
    \hline
  104861m  & 103518m  &         & 168924m  &         &161650m \\
\hline
\end{tabularx}
\end{table}


\clearpage

\section{问题二的建模与求解}

\subsection{问题二分析}
问题二在问题一的基础上，由单点往返扩展为\textbf{异构无人机多点、多架次联合调度}。其主要难点在于货箱组批、服务区访问顺序、无人机选择、架次开始时间以及共享电池周转相互耦合，若直接对所有决策变量同时搜索，将产生较大的组合规模。

因此，本问可采用\textbf{“可行架次生成 + 资源时序调度”}的分层求解思路。首先依据问题一得到的载荷与能耗计算方法，生成满足重量、体积、飞行能量和返航安全要求的候选运输架次，并确定每个架次可访问的服务区及其访问顺序；随后将这些候选架次分配给具体无人机和共享电池，并确定各架次开始时刻。调度过程中重点处理医疗物资、首批保障物资等具有较强时效要求的货箱，同时保证无人机任务占用、电池充电及再次使用之间不存在时间冲突。最终以\textbf{配送及时性、任务完成时间、总能耗和架次数}作为主要评价指标，形成兼顾救援效率与资源利用率的运输调度方案。


\section{问题三的建模与求解}
\subsection{问题三分析}
问题三的关键是将问题二中的运输调度进一步与\textbf{连续通信保障}相结合。山区环境中，运输无人机与固定网关之间的通信状态会随位置、高度和地形遮挡不断变化，因此不能仅判断服务区处是否具有通信条件，而应对无人机的整个飞行轨迹进行通信可用性检测。

首先根据运输无人机在爬升、巡航、下降和投送阶段的空间轨迹，结合DEM地形进行视距遮挡判断，并利用链路预算计算无人机与固定网关之间的通信状态。对于无法直连的轨迹区段，引入中继无人机提供通信保障。由于中继悬停位置属于连续空间变量，直接联合优化难度较大，因此可在满足地形和高度限制的区域内构造\textbf{候选中继位置集合}，预先计算各候选位置对不同运输轨迹和时间区段的通信覆盖关系，将连续空间问题转化为有限候选点上的时空覆盖问题。

在此基础上，将运输架次、中继位置、中继服务时段和能源组件进行联合调度，保证每一时刻运输无人机至少能够通过固定网关直连或中继链路保持通信。故本问的核心突破口在于\textbf{先建立轨迹级通信可行性矩阵，再进行运输--中继协同调度}，从而在保证通信连续性的同时，兼顾任务完成时间、运输与中继能耗以及设备使用规模。

\section{问题四的建模与求解}
\subsection{问题四分析}
问题四是在问题三既定联合调度方案的基础上研究\textbf{任务分区与资源独立配置}问题。由于题目要求保持原有货箱组批、访问顺序、运输任务和通信保障关系不变，因此本问并不重新规划运输路线，而是需要识别哪些服务区之间因共享同一运输架次或通信任务而具有不可分割的关联关系。

为此，可首先构建\textbf{服务区任务关联图}：将各服务区视为节点，若多个服务区出现在同一运输架次中，则在其之间建立关联，并将相互关联的服务区合并为不可拆分的任务单元。在此基础上，分别对这些任务单元进行两分区和三分区，使各组在任务量和资源需求上尽可能均衡。

完成分区后，需要根据每个任务组中的运输与中继任务时间区间，计算运输无人机、共享电池、中继无人机和能源组件的\textbf{最大并发占用量}，而不能简单累加任务总数，以此确定各分区独立运行时真正需要配置的最小资源数量。最终从资源配置规模、资源冗余、工作量均衡程度及现有库存缺口等方面比较两种分区方案，并识别导致资源不足的关键任务和时间段。

\section{模型评价}
\subsection{模型的优缺点}
\subsection{模型的改进方向}


%参考文献   手工录入
%\begin{thebibliography}{9}%宽度9
% \bibitem{bib:one} ....
% \bibitem{bib:two} ....
%\end{thebibliography}

%采用bibtex方案
\cite{mittelbach_latex_2004,wright_latex3_2009,beeton_unicode_2008,vieth_experiences_2009}

\bibliographystyle{gmcm}
\bibliography{reference}


\clearpage
%附录
\begin{appendices}
%\setcounter{page}{1} %如果需要可以自行重置页码。
\section{MATLAB 源程序}
\renewcommand{\thesubsection}{A\thinskip.\thinskip\arabic{subsection}}
\subsection{第1问程序}
\vspace{-2ex}

\begin{Matlab}{code.m}
clear all
kk=2;
[mdd,ndd]=size(dd);
while ~isempty(V)
    [tmpd,j]=min(W(i,V));
    tmpj=V(j);
    for k=2:ndd
        [tmp1,jj]=min(dd(1,k)+W(dd(2,k),V));
        tmp2=V(jj);
        tt(k-1,:)=[tmp1,tmp2,jj];
    end
    tmp=[tmpd,tmpj,j;tt];
    [tmp3,tmp4]=min(tmp(:,1));
    if tmp3==tmpd,
        ss(1:2,kk)=[i;tmp(tmp4,2)];
    else
        tmp5=find(ss(:,tmp4)~=0);
        tmp6=length(tmp5);
        if dd(2,tmp4)==ss(tmp6,tmp4)
            ss(1:tmp6+1,kk)=[ss(tmp5,tmp4);tmp(tmp4,2)];
        else, ss(1:3,kk)=[i;dd(2,tmp4);tmp(tmp4,2)];
        end
    end
    dd=[dd,[tmp3;tmp(tmp4,2)]];
    V(tmp(tmp4,3))=[];
    [mdd,ndd]=size(dd);kk=kk+1;
end;
S=ss; D=dd(1,:);
\end{Matlab}
\vspace{2ex}

\clearpage
\section{Python 源程序}
\renewcommand{\thesubsection}{B\thinskip.\thinskip\arabic{subsection}}
\subsection{第2问程序}
\vspace{-2ex}
\begin{Python}{mip1.py}
# This example formulates and solves the following simple MIP model:
#  maximize
#        x +   y + 2 z
#  subject to
#        x + 2 y + 3 z <= 4
#        x +   y       >= 1
#        x, y, z binary

# import gurobipy as gp
from gurobipy import * #GRB
try:
    # Create a new model
    m = Model("mip1")
    # Create variables
    x = m.addVar(vtype=GRB.BINARY, name="x")
    y = m.addVar(vtype=GRB.BINARY, name="y")
    z = m.addVar(vtype=GRB.BINARY, name="z")
    # Set objective
    m.setObjective(x + y + 2 * z, GRB.MAXIMIZE)
    # Add constraint: x + 2 y + 3 z <= 4
    m.addConstr(x + 2 * y + 3 * z <= 4, "c0")
    # Add constraint: x + y >= 1
    m.addConstr(x + y >= 1, "c1")
    # Optimize model
    m.optimize()
    for v in m.getVars():
        print('%s %g' % (v.varName, v.x))
    print('Obj: %g' % m.objVal)

except GurobiError as e:
    print('Error code ' + str(e.errno) + ': ' + str(e))

except AttributeError:
    print('Encountered an attribute error')
\end{Python}

\end{appendices}


\end{document}
